\chapter{Estabilidad}
\noindent
Tendremos un punto en el espacio sobre el que actúan una serie de fuerzas $F_1,\ldots,F_n$. La condición para que la partícula no se mueva es que $\sum\limits_{i=1}^{n}F_i = 0$. Este sistema de ecuaciones puede reescribirse mediante el potencial, suponiendo que todas las fuerzas son conservativas tenemos entonces que para cada fuerza $F_i$ ha de existir $V_i$ de forma que:
\begin{equation*}
    F_i = -\nabla V_i = \left(\begin{array}{ccc}
        -\frac{\partial V}{\partial x_1} \\
        \vdots \\
        -\frac{\partial V}{\partial x_d} 
    \end{array}\right), \qquad \forall i \in \{1,\ldots,n\}
\end{equation*}
Así, tenemos entonces:
\begin{equation*}
    0 = \sum_{i=1}^{n}\nabla V_i = \nabla\left(\sum_{i=1}^{n}V_i\right)
\end{equation*}
Con lo que definiendo $V = \sum\limits_{i=1}^{n}V_i$, decir que el sistema esté en equilibrio es equivalente a decir que el punto es un punto crítico del potencial $V$.

\begin{ejemplo}
    Tenemos un muelle atado al techo del que cuelga una masa $m$. La posición de la masa será función del tiempo $y(t)$. Sobre la masa actúan dos fuerzas, el peso $F_1 = -mg$ y la fuerza recuperadora del muelle $F_2$, que depende de la constante de elasticidad $k$ del muelle mediante la Ley de Hooke: $F_2 = -ky$. 
    
    \begin{figure}[H]
        \centering
        \begin{tikzpicture}[thick]

            % Techo
            \fill[gray!40] (-2,0) rectangle (2,0.3);

            % Punto de anclaje
            \fill (0,0.15) circle (2pt);

            % Muelle más corto
            \draw[
                decorate,
                decoration={coil,aspect=0.4,segment length=4pt,amplitude=4pt}
            ] (0,0.15) -- (0,-2);

            % Masa
            \draw[fill=blue!20] (-0.6,-3) rectangle (0.6,-2);

            % Etiqueta de la masa
            \node at (1,-2.5) {$m$};

            % Fuerza hacia abajo
            \draw[-{Latex[length=3mm]}, thick]
                (0.4,-2.5) -- (0.4,-4)
                node[right] {$F_1$};

            % Fuerza hacia arriba
            \draw[-{Latex[length=3mm]}, thick]
                (0.4,-2.5) -- (0.4,-1)
                node[right] {$F_2$};

            % Fuerza hacia arriba
            \draw[-{Latex[length=3mm]}, thick]
                (-1.4,-3.5) -- (-1.4,-0.5)
                node[right] {$y$};

            \fill(0.4,-2.5) circle(2pt) [] node{};
        \end{tikzpicture}
    \end{figure}

    \noindent
    Aplicando la ecuación anterior suponiendo que nos encontramos en una posición $y_\ast$ de equilibrio:
    \begin{equation*}
        F_1 + F_2 = -mg -ky_\ast = 0 \quad\Longrightarrow\quad y_\ast = -\frac{mg}{k}
    \end{equation*}
    Si queremos hacer la fórmula anterior con los potenciales, tenemos:
    \begin{equation*}
        V_1(y) = mgy, \qquad V_2(y) = \frac{ky^2}{2}
    \end{equation*}
    Por lo que:
    \begin{equation*}
        V(y) = mgy + \frac{ky^2}{2}
    \end{equation*}
    Es una parábola, por lo que tiene un único punto crítico, que tiene que ser el punto de equilibrio $y_\ast$.
\end{ejemplo}

\begin{ejemplo}
    Un potencial que presenta más de un punto crítico puede ser como ejemplo:
    \begin{equation*}
        V(\theta) = k(1-\cos\theta)
    \end{equation*}
    Que representa el potencial asociado a un péndulo. Tenemos una función $V:\mathbb{R}\to\mathbb{R}$ $2\pi-$periódica.
\end{ejemplo}

\section{Soluciones estables}
\noindent
El marco teórico en el que trabajaremos será en el que tenemos una ecuación diferencial:
\begin{equation*}
    \dot{x} = X(t,x)
\end{equation*}
donde $X:D\subset\mathbb{R}\times\mathbb{R}^d\to\mathbb{R}^d$ es un campo continuo de dominio conexo y abierto, supondremos que hay unicidad para toda condición inicial, por lo que podremos usar la notación $x(t;t_0,x_0)$.

\begin{definicion}
    Diremos que una solución $x(t;t_0,x_0)$ es \underline{estable} (en el sentido de Lyapunov) si está definida en $\left[t_0,+\infty\right[$ y:
        \begin{equation*}
            \forall \varepsilon>0~~\exists \delta>0~:~\|x_0-\widetilde{x_0}\| < \delta \quad \text{entonces:}
        \end{equation*}
        \begin{itemize}
            \item $x(t;t_0,\widetilde{x_0})$ está definida en $\left[t_0,+\infty\right[$
            \item $\|x(t;t_0,x_0)-x(t;t_0,\widetilde{x_0})\| < \varepsilon\quad \forall t\geq t_0$
        \end{itemize}
\end{definicion}

\begin{ejercicio}
    Esta definición es independiente de $t_0$. Es decir, $x(t;t_0,x_0)$ es estable si y solo si lo es $x(t;t_1,x_1)$. Este ejercicio se sigue del Teorema de dependencia continua respecto a condiciones iniciales.
\end{ejercicio}

\begin{ejemplo}
    Familiaricémonos con este concepto.
    \begin{enumerate}
        \item Consideramos la ecuación:
            \begin{equation*}
                \dot{x} = 1
            \end{equation*}
            Cuyas soluciones son $x(t) = t+c$. Todas estas soluciones son estables. Observemos que:
            \begin{equation*}
                x(t;t_0,x_0) = t-t_0+x_0, \quad t\in \mathbb{R}
            \end{equation*}
            Fijada esta solución, para $\delta=\varepsilon$ se tiene claramente la estabilidad de $x(t;t_0,x_0)$.
        \item Ahora para la ecuación:
            \begin{equation*}
                \dot{x} = -x+1+t
            \end{equation*}
            Vemos que $x(t;0,0) = t$ sigue siendo una solución. ¿Será estable? Se trata de una ecuación lineal completa, tenemos ya suna solución particular, por lo que la solución general será:
            \begin{equation*}
                x(t;0,x_0) = t + x_0e^{-t}, \quad t\in \mathbb{R}
            \end{equation*}
            Vemos que la solución $x(t;0,0)$ sigue siendo estable, pues si $|x_0|<\delta$, para $\delta = \varepsilon$ tenemos que:
            \begin{equation*}
                |x(t;0,x_0)-x(t;0,0)| = |x_0|e^{-t} \leq |x_0| < \delta = \varepsilon, \quad \forall t\geq 0
            \end{equation*}
        \item Si consideramos finalmente la ecuación
            \begin{equation*}
                \dot{x} = x+1-t
            \end{equation*}
            Vemos que $x(t;0,0)=t$ sigue siendo una solución, que ahora no será estable. Tenemos:
            \begin{equation*}
                x(t;0,x_0) = t + x_0e^{t}, \quad t\in \mathbb{R}
            \end{equation*}
            Por reducción al absurdo, si fuera estable, para $\varepsilon=1$ podríamos encontrar $\delta>0$ de forma que si $|x_0|<\delta$ tendríamos entonces que:
            \begin{equation*}
                |x(t;0,x_0)-x(t;0,0)| = |x_0| e^t < 1 \quad \forall t\geq 0
            \end{equation*}
            Lo cual es imposible para valores $x_0 \neq 0$, pues $|x_0|e^t$ no es una función acotada.
    \end{enumerate}
\end{ejemplo}

\begin{definicion}
    Diremos que una solución $x(t;t_0,x_0)$ es \underline{asintóticamente estable} si:
    \begin{itemize}
        \item $x(t;t_0,x_0)$ es estable.
        \item $\exists \rho>0$ tal que si $\|x_0-\widetilde{x_0}\|<\rho$ entonces:
            \begin{equation*}
                \lim_{t\to+\infty}\|x(t;t_0,x_0)-x(t;t_0,\widetilde{x_0})\| = 0
            \end{equation*}
    \end{itemize}
\end{definicion}

\noindent
Además, es importante destacar que la segunda condición no implica la primera.

\begin{ejemplo}
    En el ejemplo anterior de $\dot{x} = 1$ teníamos que las soluciones eran estables pero no asintóticamente estables y en la ecaución $\dot{x} = -x+1+t$ la solución $x(t) = t$ es asintóticamente estable.
\end{ejemplo}~\\

\noindent
En lo que queda de sección, vamos a considerar solo \underline{ecuaciones autónomas}, donde el campo $X$ no depende de $t$:
\begin{equation*}
    \dot{x} = X(x)
\end{equation*}
Así, tendremos un campo $X:\Omega\subset\mathbb{R}^d\to\mathbb{R}^d$ continuo con $\Omega$ un conjunto abierto y conexo, suponiendo además unicidad de soluciones para toda condición inicial. Podemos ver esta teoría dentro de la anterior considerando $D=\mathbb{R}\times\Omega$ y $t_0=0$.\\

\noindent
Para el estudio de estas ecuaciones estaremos interesados al principio en buscar \underline{soluciones constantes} o \underline{equilibrios}, es decir, en buscar puntos $x_\ast \in \Omega$ de forma que $X(x_\ast)=0$. Encontrado un punto como este, la función:
\begin{equation*}
    x(t) = x_\ast, \quad t\in \mathbb{R}
\end{equation*}
Es una solución constante de la ecuación.

\begin{ejemplo}
    \textbf{Ecuación logística.} Se trata de una ecuación que trata de modelar el creciemiento de algo, ya sean poblaciones u organismos. Viene dada por:
    \begin{equation*}
        \dot{x} = k\left(1-\frac{x}{M}\right)x
    \end{equation*}
    para $k,M>0$. Viene motivada por la ecuación de Mathus:
    \begin{equation*}
        \dot{x} = \lm x
    \end{equation*}
    $k$ mide lo rápido que crece la población y $M$ mide la ``capacidad de carga'', lo que puede soportar el hábitat.\\

    \noindent
    Pasando a teoría de ecuaciones diferenciales, es una ecuación de primer orden autónoma con campo $X:\mathbb{R}\to\mathbb{R}$ dado por:
    \begin{equation*}
        X(x) = xk\left(1-\frac{x}{M}\right)
    \end{equation*}
    Los equilibrios de esta ecuación son $0$ y $M$. Es una ecuación de variables separadas, que tras resolverla obtenemos tres tipos de soluciones, en relación de las situaciones $x<0$, $0<x<M$ o $x>M$.\\

    \noindent
    Al graficarlas se observa que la solución constantemente igual a $M$ es asintóticamente estable y que la sollución constantemente igual a $0$ no es estable. Hágase de forma analítica.
\end{ejemplo}

\subsection{Sistemas lineales de coeficientes constantes}
\noindent
Consideraremos sistemas del tipo:
\begin{equation*}
    \dot{x} = Ax
\end{equation*}
para $x\in \mathbb{R}^d$ y $A\in \mathbb{R}^{d\times d}$. Tenemos espacio $\Omega=\mathbb{R}^d$. En este caso las soluciones constantes se corresponden a las constantes $\ker A$. Usualmente solo se tiene la solución constantemente igual a cero, por lo que centraremos nuestro estudio en ver bajo qué casos esta solución es estable o no.\\

\noindent
Lo que hacíamos en Ecuaciones Diferenciales I para buscar soluciones de esta ecuación era buscar valores propios, pues una vez tenemos:
\begin{equation*}
    Av = \lm v
\end{equation*}
Podemos considerar la solución $x(t) = e^{\lm t}v, \quad t\in \mathbb{R}$, ya que:
\begin{equation*}
    \dot{x}(t) = \lm e^{\lm t}v = e^{\lm t} Av = Ax(t) \qquad \forall t\in \mathbb{R}
\end{equation*}
Estas soluciones son curvas paramétricas que describen semirrectas de origen el cero. Estas semirrectas son repulsivas de $0$ cuando $\lm>0$ y atractivas a $0$ cuando $\lm<0$.\\

\noindent
Así, cuando $\lm>0$ está claro que la solución constantemente igual a $0$ es inestable, pues basta tomar una condición inicial en la dirección de $v$ para obtener una solución que escapa de $0$. Así, para que la solución sea estable, hay que exigir que todos los valores propios de $A$ sean negativos.\\

\noindent
Sin embargo, lo discutido solo se cumple para matrices cuyos valores propios sean reales. ¿Qué ocurre en este caso? Supongamos que tenemos $\lm\in \mathbb{C}\setminus \mathbb{R}$ valor propio de la matriz $A$, con lo que:
\begin{equation*}
    A w = \lm w
\end{equation*}
para $w = u + iv$, $v\neq 0$. Como $A$ es una matriz real, sabemos que los valores propios conjugados van en pares, pues si $\lm\in \mathbb{C}\setminus \mathbb{R}$ es un valor propio entonces:
\begin{equation*}
    A\overline{w} = \overline{Aw} = \overline{\lm w} = \overline{\lm}\overline{w}
\end{equation*}
Los vectores propios reales generan rectas invariantes por una transformación lineal. Por otra part, los vectores propios complejos generan planos invariantes, como en el caso de una rotación.\\

\noindent
Si $\lm = a + ib$ y $w= u+iv$ con $b,v\neq 0$ de forma que:
\begin{equation*}
    Aw = \lm w
\end{equation*}
entonces podemos considerar la solución:
\begin{align*}
    x(t) &= e^{\lm t}w = e^{at} e^{ibt}(u + iv) = e^{at}(\cos (bt) + i\sen (bt))(u+iv) \\
         &= e^{at}[u\cos (bt)-v\sen(bt) + i (u\sen(bt) + v\cos(bt))]
\end{align*}
De donde tomando partes real e imaginaria obtenemos las soluciones:
\begin{align*}
    e^{at}(u\cos(bt) - v\sen(bt)) \\
    e^{at}(u\sen(bt) + v\cos(bt))
\end{align*}
Vemos que $a$ es un factor de dilatación del giro, y la expresión restante es una elipse de ejes $u$ y $v$. Cuando $a=0$ tenemos elipses concéntricas y cuando $a\neq0$ obtenemos espirales elípticas.\\

\noindent
Así, en el caso complejo, la estabilidad de la solución constantement igual a cero depende de $a = Re(\lm)$.
